% language: swedish
% figure: fig1 0.055 0.391 0.375 0.465
\section{Föreläsning 5/9}

\begin{theorem}[Cauchy-Riemanns ekvation (Sats 2.13)]
Låt $G \subseteq \mathbb{C}$, $f : G \to \mathbb{C}$, $u = \operatorname{Re} f$, $v = \operatorname{Im} f$
\begin{enumerate}[label=\alph*)]
  \item Om $f$ \rattat{är}{hade} holomorf i $G$ så löser $u, v$ CRs ekvationer:
  \[ \begin{cases} \dfrac{du}{dx} = \dfrac{dv}{dy} \\[2ex] \dfrac{du}{dy} = -\dfrac{dv}{dx} \end{cases} \qquad \text{alt.} \quad \frac{df}{dx} = -i \frac{df}{dy} \quad (= f') \]
  \item Om $f \in C^1$ (har kontinuerliga partiella derivator) och $u, v$ löser CRs ekvationer, då är $f$ holomorf i $G$.
\end{enumerate}
\end{theorem}

\textbf{Geometriskt:}
\notefigure{fig1}{}
\textcolor{magenta!80!black}{Finns en youtube video på kurshemsidan som förklarar hur holomorf funkar.}

\begin{example}
Antag att $f$ är holomorf i $\mathbb{C}$ och $\operatorname{Re} f = x^2 - y^2$. Bestäm $\operatorname{Im} f$.

\textbf{Lösning:} Enligt CRs ekv.\ löser $u = \operatorname{Re} f$ och $v = \operatorname{Im} f$ ekv.
\[ \begin{cases} \dfrac{du}{dx} = \dfrac{dv}{dy} \\[2ex] \dfrac{du}{dy} = -\dfrac{dv}{dx} \end{cases} \]
får $\rattat{\dfrac{du}{dx}}{\dfrac{du}{dy}} = 2x$ \quad $\dfrac{dv}{dx} = 2y$ \quad $v = 2xy + g(x)$, \ $\dfrac{dv}{dx} = 2y$

$\Rightarrow g'(x) = 0 \Rightarrow v = 2xy + C$ \qquad $(f(z) = z^2 + iC)$
\end{example}

\begin{theorem}[(2.3 PB, 12.64 Ad)]
Om $f \in C^1$, då är $f$ \emph{differentierbar}, dvs
\begin{align*}
  f(x_0 + t, y_0 + s) - f(x_0, y_0) &= \\
  &= \frac{df}{dx}(x, y)t + \frac{df}{dy}(x_0, y_0)s + O\left(\sqrt{t^2 + s^2}\right)
\end{align*}
\end{theorem}
